
\documentclass[11pt]{article}

\usepackage[utf8]{inputenc}
\usepackage[T1]{fontenc}
\usepackage{lmodern}
\usepackage[empty]{fullpage}
\usepackage{titlesec}
\usepackage[usenames,dvipsnames]{color}
\usepackage{enumitem}
\usepackage[hidelinks]{hyperref}
\usepackage{fancyhdr}
\usepackage[english,portuguese]{babel}
\usepackage{tabularx}

\pagestyle{fancy}
\fancyhf{} 
\fancyfoot{}
\renewcommand{\headrulewidth}{0pt}
\renewcommand{\footrulewidth}{0pt}

\addtolength{\oddsidemargin}{-0.5in}
\addtolength{\evensidemargin}{-0.5in}
\addtolength{\textwidth}{1.0in}
\addtolength{\topmargin}{-.6in}
\addtolength{\textheight}{1.2in}

\urlstyle{same}
\raggedbottom
\raggedright
\setlength{\tabcolsep}{0in}

\titleformat{\section}{
  \vspace{-4pt}\scshape\raggedright\large\bfseries
}{}{0em}{}[\color{black}\titlerule \vspace{-5pt}]

\newcommand{\resumeItem}[1]{
  \item\small{
    {#1 \vspace{-2pt}}
  }
}

\newcommand{\resumeSubheading}[4]{
  \vspace{-2pt}\item
    \begin{tabular*}{1.0\textwidth}[t]{l@{\extracolsep{\fill}}r}
      \textbf{#1} & \textbf{\small #2} \\
      \textit{\small#3} & \textit{\small #4} \\
    \end{tabular*}\vspace{-7pt}
}

\newcommand{\resumeProjectHeading}[2]{
    \item
    \begin{tabular*}{1.0\textwidth}{l@{\extracolsep{\fill}}r}
      \small#1 & \textbf{\small #2} \\
    \end{tabular*}\vspace{-7pt}
}

\renewcommand\labelitemi{$\vcenter{\hbox{\tiny$\bullet$}}$}
\renewcommand\labelitemii{$\vcenter{\hbox{\tiny$\bullet$}}$}

\newcommand{\resumeSubHeadingListStart}{\begin{itemize}[leftmargin=0.0in, label={}]}
\newcommand{\resumeSubHeadingListEnd}{\end{itemize}}
\newcommand{\resumeItemListStart}{\begin{itemize}[leftmargin=0.15in]}
\newcommand{\resumeItemListEnd}{\end{itemize}\vspace{-5pt}}

\begin{document}

\begin{center}
    {\Huge \textbf{Mateus Amaral da Silva}} \\ \vspace{5pt}
    \small (31) 99136-9726 \quad
    \href{mailto:silvaamaralmateus@gmail.com}{\underline{silvaamaralmateus@gmail.com}} \quad
    \href{https://www.linkedin.com/in/mateus-amaral-da-silva}{\underline{LinkedIn}} \quad
    \href{https://github.com/MateusAmaralDaSilva}{\underline{GitHub}} \quad
    \href{https://mateus-amaral-site.vercel.app}{\underline{Portfólio}}
\end{center}

\section{Resumo Profissional}
 \small{Cientista da Computação focado em IA aplicada a ambientes sensíveis, com experiência prática em aprendizado profundo e aprendizado federado, além de sólido conhecimento em segurança e desenvolvimento de aplicações web (fullstack). Atuo desenvolvendo modelos para detecção de câncer de colo uterino e conduzindo análises de performance para estes sistemas. Tenho experiência em auxiliar e ensinar membros iniciantes em IA/ML, e interesse em aplicar IA preditiva e generativa em produção com foco em eficiência, impacto, escalabilidade e segurança do usuário.}

\section{Experiência Profissional}
  \resumeSubHeadingListStart

    \resumeSubheading
      {Cientista de Dados}{Abr 2026 -- Jul 2026}
      {Guidance IA (Empresa de consultoria de IA)}{Remoto}
      \resumeItemListStart
        \resumeItem{Trabalhei com sistemas multiagentes para uma empresa de advocacia, criando agentes confiáveis e seguros para automação e criação de contratos jurídicos.}
        \resumeItem{Desenvolvi um sistema de sumarização de tickets multiagente para a empresa Ânima em um ambiente de grande volume de dados, reduzindo o custo de tokens e elevando a precisão geral do modelo.}
      \resumeItemListEnd

    \resumeSubheading
      {Pesquisador -- Projeto CI-IA Saúde (Detecção de Câncer de Colo Uterino)}{Abr 2025 -- Presente}
      {Universidade Federal de Minas Gerais (UFMG)}{Belo Horizonte, MG}
      \resumeItemListStart
        \resumeItem{Implementei modelos de aprendizado federado na plataforma do laboratório (Flautim) voltados à detecção de câncer de colo uterino, alcançando acurácia superior a 95\% mesmo em cenários de dados altamente desbalanceados.}
        \resumeItem{Apliquei métodos avançados de avaliação de modelos e Explainable AI (XAI) para interpretação rigorosa e validação das decisões do sistema.}
        \resumeItem{Atuei como representante da equipe de engenharia de dados e modelos, auxiliando na organização técnica e na capacitação de novos membros em machine learning e aprendizado federado.}
      \resumeItemListEnd

    \resumeSubheading
      {Desenvolvedor e Gestor de Produção}{Fev 2022 -- Dez 2022}
      {IFMG | Solveware}{Minas Gerais}
      \resumeItemListStart
        \resumeItem{Liderei a comunicação direta com clientes para levantamento de requisitos de um sistema de gerenciamento de condomínios.}
        \resumeItem{Desenvolvi a modelagem relacional do banco de dados e a implementação do front-end da aplicação.}
      \resumeItemListEnd

  \resumeSubHeadingListEnd

\section{Educação}
  \resumeSubHeadingListStart
    \resumeSubheading
      {Universidade Federal de Minas Gerais (UFMG)}{Belo Horizonte, MG}
      {Bacharelado em Ciência da Computação (7º Semestre -- Previsão: Junho de 2028)}{Mar 2023 -- Presente}
    \resumeSubheading
      {Instituto Federal de Educação, Ciência e Tecnologia de Minas Gerais (IFMG)}{Minas Gerais}
      {Ensino Técnico em Informática}{2020 -- 2022}
  \resumeSubHeadingListEnd

\section{Principais Projetos Pessoais}
    \resumeSubHeadingListStart
      \resumeProjectHeading
        {\textbf{Sistema Multiagente para NPC's Imersivos em Jogos RPG}}{GitHub / Site}
        \resumeItemListStart
          \resumeItem{Desenvolvimento de agentes inteligentes para jogos estilo RPG com foco em comportamento imersivo, acompanhados de métricas detalhadas de performance e interpretabilidade.}
        \resumeItemListEnd
      \resumeProjectHeading
        {\textbf{Blog Pessoal e Portfólio Full-Stack}}{GitHub / Site}
        \resumeItemListStart
          \resumeItem{Aplicação full-stack desenvolvida e hospedada seguindo rigorosas boas práticas de engenharia de software, Next.js e design moderno.}
        \resumeItemListEnd
    \resumeSubHeadingListEnd

\section{Competências Técnicas}
 \begin{itemize}[leftmargin=0.15in, label={}]
    \small{\item{
     \textbf{Linguagens:}{ Python, SQL, C++, Java, HTML, CSS, TypeScript, JavaScript} \\
     \textbf{Libs / Frameworks:}{ PyTorch, TensorFlow, Scikit-learn, FastAI, FastAPI, Next.js, React} \\
     \textbf{MLOps \& DevOps:}{ Flautim, Flower, Docker, Git} \\
     \textbf{Bancos de Dados:}{ MySQL, PostgreSQL} \\
     \textbf{Conhecimentos:}{ Aprendizado centralizado e federado, privacidade e segurança de dados, métricas avançadas de validação, data augmentation, feature engineering, desenvolvimento full-stack e boas práticas de programação.}
    }}
 \end{itemize}

\section{Idiomas}
 \begin{itemize}[leftmargin=0.15in, label={}]
    \small{\item{
     \textbf{Português:} Nativo \quad | \quad
     \textbf{Inglês:} Avançado \quad | \quad
     \textbf{Espanhol:} Intermediário \quad | \quad
     \textbf{Libras:} Básico
    }}
 \end{itemize}

\end{document}
